\chapter{Conclusion}\label{ch:conclusion}

On a shared budget of 20,000 AG News training articles, the simple sparse linear baseline (\textit{Word TF-IDF LR}) is the clear winner in clean accuracy, typo resistance, and speed. It reached 90.05\% accuracy on clean articles and dipped by less than one percentage point (to 89.14\%) when 20\% of eligible words contained swapped interior letters, fitting in just 13 seconds on standard CPU hardware. Both convolutional networks fell noticeably faster under noise (\textit{Word CNN} by 2.31 points and \textit{Subword CNN} by 1.86 points) while taking 17 to 50 times longer to train.

The core explanation lies in how each model weighs textual clues across an article. Because additive linear scoring sums dozens of surviving words and word pairs, losing a single misspelled token barely disturbs the overall decision. Global max pooling inside CNNs, on the other hand, funnels each filter through its single highest-scoring word window, triggering more than twice as much article-level decision churn, concentrating losses in \emph{Business}, and spilling ambiguous stories into \emph{World}.

Training on a 50/50 mix of clean and 10\%-swapped articles gave \textit{Subword CNN} a small average edge under heavy noise (0.31 to 0.46 points across random typo sets), though run-to-run variation and a slight clean-text drop on unseen articles limit its consistency. Extending the word-only CNN to the same 16-epoch decay schedule also closed nearly all of \textit{Subword CNN}'s augmented lead, showing that training schedule length accounted for much of the difference.

For classifying short news items under moderate spelling noise, clean-trained \textit{Word TF-IDF LR} remains the strongest, fastest, and most dependable default. Richer subword neural architectures or dynamic training noise are best reserved for tasks where much larger datasets, matched training budgets, and GPU acceleration justify the added cost.
